% ── Gift 诗歌赏析书：宏包与自定义环境 ──

\usepackage{geometry}
\geometry{a5paper, margin=2.5cm, bottom=2cm}

\usepackage{fontspec}
\usepackage{graphicx}

% ── CJK 字体 ──
% fontset=none 后手动定义所有 CJK 字体族
\setCJKmainfont{SourceHanSerifSC-Regular.otf}[Path=./assets/fonts/]
\setCJKfamilyfont{zhsong}{SourceHanSerifSC-Regular.otf}[Path=./assets/fonts/]
\setCJKfamilyfont{zhkai}{LXGWWenKai-Regular.ttf}[
    Path=./assets/fonts/,
    BoldFont=LXGWWenKai-Medium.ttf,
]
\newcommand{\songti}{\CJKfamily{zhsong}}
\newcommand{\kaishu}{\CJKfamily{zhkai}}

\usepackage{xcolor}
\usepackage{titlesec}
\usepackage{fancyhdr}
\usepackage{hyperref}

% ── 颜色 ──────────────────────────
\definecolor{pagebg}{HTML}{f9f4ef}
\definecolor{accent}{HTML}{9e4a3e}
\definecolor{textgray}{HTML}{333333}
\definecolor{warmbrown}{HTML}{6b4c3b}
\definecolor{teal}{HTML}{4a7c7c}

\pagecolor{pagebg}

\hypersetup{
    colorlinks=true,
    linkcolor=accent,
    urlcolor=accent,
}

% ── 页码风格 ──────────────────────
\pagestyle{fancy}
\fancyhf{}
\fancyfoot[C]{\small\thepage}
\renewcommand{\headrulewidth}{0pt}

% ── 章节标题格式 ──────────────────
\titleformat{\chapter}[display]
    {\centering\Huge\kaishu}
    {}
    {0pt}
    {\vspace*{-2em}}
    []
\titlespacing{\chapter}{0pt}{40pt}{30pt}

% 带编号的 section（备用）
\titleformat{\section}[block]
    {\centering\large\kaishu}
    {\thesection}
    {0.5em}
    {}
\titlespacing{\section}{0pt}{20pt}{8pt}

% 无编号 section（诗歌标题用）
\titleformat{name=\section,numberless}[block]
    {\centering\large\kaishu}
    {}
    {0pt}
    {}
\titlespacing*{name=\section,numberless}{0pt}{20pt}{8pt}

% ── 诗歌环境 ──────────────────────
% 居中，宋体，字号略小于正文，行距 2 倍
\newenvironment{poem}
    {\par\addvspace{1.5em}
     \begingroup
     \centering
     \color{warmbrown}
     \songti\normalsize
     \setlength{\baselineskip}{2em}
     \obeylines}
    {\endgroup
     \par\addvspace{1.5em}}

% ── 赏析环境 ──────────────────────
% 宋体，字号 11bp（略小于正文），行距 1.6 倍，首行缩进
\newenvironment{appreciation}
    {\par\addvspace{0.5em}
     \songti\zihao{5}
     \setlength{\baselineskip}{1.75em}
     \setlength{\parindent}{2em}
     \setlength{\parskip}{0.5em}}
    {\par\addvspace{0.5em}}

% ── 朝代导言环境 ──────────────────
% 宋体，正常字号，首行缩进，用于每章开篇介绍
\newenvironment{intro}
    {\par\addvspace{0.5em}
     \songti\normalsize
     \setlength{\baselineskip}{1.75em}
     \setlength{\parindent}{2em}
     \setlength{\parskip}{0.5em}}
    {\par\addvspace{0.5em}}

% ── 注释环境 ──────────────────────
% 字号更小（\footnotesize），楷体，靛青色，用于诗歌正文后的注释块
\newenvironment{annotations}
    {\par\addvspace{1em}
     \color{teal}
     \footnotesize\kaishu
     \setlength{\parindent}{0em}
     \setlength{\parskip}{0.3em}}
    {\par\addvspace{1em}}

\newcommand{\notetext}[2]{%
    \noindent #1.\enspace #2\par}

% ── 元信息标记 ────────────────────
% 每首诗前自动显示标题、作者、朝代（无序号，居中，加入目录）
\newcommand{\poemheading}[3]{%
    \phantomsection
    \section*{#1}%
    \addcontentsline{toc}{section}{#1}%
    \begin{center}
    \small\kaishu #3 · #2
    \end{center}
    \vspace{-0.5em}
}
